\documentclass[10pt,a4paper]{amsart}
\usepackage[margin=19mm]{geometry}
\usepackage{amsmath,amssymb,mathtools}
\usepackage[scr=rsfs,cal=euler]{mathalfa}
\usepackage{xfrac}
\usepackage[unicode,hidelinks,bookmarks=false]{hyperref}
\newcommand{\ie}{i.e.\ }
\newcommand{\cf}{cf.\ }
\newcommand{\resp}{respectively\ }
\newcommand{\Subsection}{\subsection}
\providecommand{\nobreakdash}{\nobreak\hbox{-}}
\setlength{\emergencystretch}{2em}
\multlinegap0pt
\usepackage[shortlabels]{enumitem}
\usepackage{bm}
\usepackage{multirow}
\usepackage{amssymb}
\usepackage{xcolor}
\usepackage{booktabs}
\usepackage{algorithm}\usepackage{algpseudocode}
\newtheorem{theorem}{Theorem}
\newtheorem{lemma}{Lemma}
\newtheorem{proposition}{Proposition}
\title{Geometry-dependent Ekman layer approximations on curved domains: $L^\infty$ convergence}
\author{Yifei Jia, Yi Du, Lihui Guo}
\date{Source: arXiv:2512.18579v1 (CC BY 4.0)}
\begin{document}
\maketitle

\noindent\emph{Source and licence.} Geometry-dependent Ekman layer approximations on curved domains: $L^\infty$ convergence. Version arXiv:2512.18579v1 (CC BY 4.0), distributed by arXiv under the Creative Commons Attribution 4.0 International licence, http://creativecommons.org/licenses/by/4.0/, as recorded in the arXiv licence metadata for this exact version and shown on the abstract page https://arxiv.org/abs/2512.18579v1. Full licence notice: \texttt{licenses/arxiv-2512.18579-CC-BY-4.0.txt}.\par\smallskip

\noindent\textit{Editorial scope.} Counted unit: the article's proposition prop8 (source0.tex lines 2925--2936). The article's complete original proof of it is delivered with it (source0.tex lines 2937--3247, one \texttt{proof} environment of 311 lines). No mathematics is written, repaired or re-derived: the statement and the proof are the article's own lines, in the article's own notation, and no retained line is substituted or re-flowed.  Retained uncounted as the source-local context the proof needs: lines 310--316; lines 318--320; lines 410--474; lines 509--544; lines 520--522; lines 620--630; lines 2353--2364; lines 2389--2428; lines 2573--2612; lines 2678--2685; lines 2704--2708; lines 2809--2836.  Nothing was omitted from any retained span, so the package carries no omission record.  No retained line contains a citation, so the wrapper carries no bibliography and no bibliography entry.  No retained line contains a figure environment, an image-inclusion command or drawing code, so no image is delivered and the package declares no asset dependency.  Editorial formatting only, in addition: an AMS \texttt{amsart} preamble that restates the article's own declarations and macros used by the retained lines, namely the environments \texttt{theorem}, \texttt{lemma}, \texttt{proposition} on separate counters, together with the standard packages those lines need.  The retained environments print in this document as Theorem 1, Lemma 1, Proposition 1 in order of appearance, whereas the article prints the same items with the numbering of its own full document.  The complete native source of the article as distributed in the arXiv e-print for this exact version is bundled unchanged as \texttt{sources/191403/source0.tex}.
\begin{equation}\label{1.11}
	\begin{cases}
		\partial_t \boldsymbol{u}^\varepsilon-\nu\varepsilon \Delta \boldsymbol{u}^\varepsilon+(\boldsymbol{u}^\varepsilon \cdot \nabla)\boldsymbol{u}^\varepsilon+\varepsilon^{-1}\boldsymbol{R}\boldsymbol{u}^\varepsilon+\varepsilon^{-1}\nabla p^\varepsilon=0\,,\\
		\nabla\cdot \boldsymbol{u}^\varepsilon=0\,,\\
		\boldsymbol{u}^\varepsilon|_{\partial\Omega}=0\,,\quad \boldsymbol{u}^\varepsilon|_{t=0}=\boldsymbol{u}_0^\varepsilon\,,
	\end{cases}  (t,x)\in \mathbb{R}^+\times \Omega\,,
\end{equation}

\begin{equation}\label{1.12}
	\Omega=\mathbb{R}^2 \times [B(x, y), B(x, y) + 2]\,,\qquad  B(x,y)\in C^\infty(\mathbb{R}^2)\,.
\end{equation}

\begin{theorem}\label{thapp}

As $\varepsilon\rightarrow 0^+$, assume $\boldsymbol{u}_0^\varepsilon$ strongly converges to ${\boldsymbol{u}}_0 (x,y)= ({\boldsymbol{u}}_{0,h}, 0)$ in $L^2(\Omega)$ and ${\boldsymbol{u}}_{0,h}\in H^6(\mathbb{R}^2)$, the boundary surface $B(x,y)\in C^\infty(\mathbb{R}^2)$ satisfies 	
	\begin{equation}\label{+16}
		\frac{\cos^2\alpha+\cos^2\beta}{\cos^2\gamma}<\tfrac{1}{8}\,,
		\quad  
		\sup\limits_{(x,y)\in \mathbb{R}^2} (1+\sqrt{\tfrac{2}{\nu}})\sqrt{\big| \boldsymbol{K}_{\mathrm{G}}\big|} <\tfrac{8}{9}\,,
	\end{equation}
where   $\boldsymbol{K}_{\mathrm{G}}$ is  the Gaussian curvature of $B(x,y)$.
Then there exists a pair of approximate solution of system \eqref{1.11}-\eqref{1.12}:
	\begin{equation}\label{1.16}
		\begin{cases}
			\boldsymbol{u}_{app}^{\varepsilon}=\sum\limits_{i=0}^{2}{\delta}^i
			 \boldsymbol u^{i}(t,x,y,z,\tilde z,\bar{z}) 
			+\boldsymbol u^{c},\\
			{p}_{app}^{\varepsilon}=\sum\limits_{i=0}^{3}{\delta}^i
			  p^{i}(t,x,y,z,\tilde z,\bar{z}) 
			+p^{c}.
		\end{cases}
	\end{equation}

	Additionally, the approximate solution satisfies:
	\begin{equation}\label{1.19}
		\begin{cases}
			\partial_{t} {\boldsymbol u}_{app}^{\varepsilon}-\nu\varepsilon\Delta{\boldsymbol u}_{app}^{\varepsilon}+({\boldsymbol u}_{app}^{\varepsilon} \cdot \nabla) {\boldsymbol u}_{app}^{\varepsilon}+\varepsilon^{-1}\boldsymbol{R} {\boldsymbol u}_{app}^{\varepsilon}
			+\varepsilon^{-1}\nabla{p}_{app}^{\varepsilon}=\boldsymbol \rho^{\varepsilon}\,,\\
			\nabla\cdot {\boldsymbol u}_{app}^{\varepsilon}=0\,,\quad
			{\boldsymbol u}_{app}^{\varepsilon}|_{\partial\Omega}=0\,,\\
			{\boldsymbol u}_{app}^{\varepsilon}|_{t=0}=\boldsymbol{u}_0^\varepsilon,
		\end{cases}
	\end{equation}
	where $\boldsymbol \rho^{\varepsilon}$ satisfies
	\begin{equation*}
		\lVert
		\boldsymbol \rho^{\varepsilon}\rVert_{{L}^{2}(\Omega)}
		\lesssim
		\varepsilon^2
		\lVert
		{\boldsymbol{u}}_{0,h}
		\rVert_{{H}^{5}(\mathbb{R}^2)}
		\big(
		\lVert
		{\boldsymbol{u}}_{0,h}
		\rVert_{{H}^{3}(\mathbb{R}^2)}
		+1\big)^2
		\notag,
	\end{equation*}	
	and
	\begin{equation*}
		\lVert
		\partial_t\boldsymbol \rho^{\varepsilon}\rVert_{{L}^{2}(\Omega)}
		\lesssim
		\varepsilon^2
		\lVert
		{\boldsymbol{u}}_{0,h}
		\rVert_{{H}^{6}(\mathbb{R}^2)}
		\big(
		\lVert
		{\boldsymbol{u}}_{0,h}
		\rVert_{{H}^{4}(\mathbb{R}^2)}
		+1
		\big)^2
		\notag.
	\end{equation*}	
\end{theorem}

\begin{theorem}\label{thm2}
	Let $\boldsymbol u^\varepsilon\in{L}^{\infty}([0,T^*);{H}^{2} (\Omega))\cap{L}^{2}([0,T^*);{H}^{3} (\Omega))$ be a family of local strong solutions of the system \eqref{1.11}-\eqref{1.12} associated with the initial data $\boldsymbol u_0^\varepsilon\in   H^{2}(\Omega)$. Under the following conditions:
	\begin{enumerate}[1.]
		\item{Well-prepared initial data:} 
		let 
		\begin{equation}\label{1.21}
			\mathop{\lim}\limits_{\varepsilon\rightarrow 0^+} \boldsymbol u^\varepsilon_0=
			({\boldsymbol u}_{0,h}, 0)
			\quad\text{with}\quad{\boldsymbol u}_{0,h}\in H^s(\mathbb{R}^2)~(s>5)\,,
		\end{equation}
		and
		\begin{equation}\label{1.26}   
		\lVert{\boldsymbol{u}}_{0,h}\rVert_{{W}^{1,\infty}(\mathbb{R}^2)}\leq\sqrt{\tfrac{\nu}{3}}\,;
		\end{equation}
		\item{The boundary $B(x,y)\in C^\infty(\mathbb{R}^2)$ satisfies:}
		\begin{align}
			|{B}(x,y)|<\tfrac14,\quad   \tfrac{\cos^2\alpha+\cos^2\beta}{\cos^2\gamma}<\tfrac{1}{8}\,,\label{1.23}\\
			\sup\limits_{(x,y)\in \mathbb{R}^2} (1+\sqrt{\tfrac{2}{\nu}})\sqrt{\big| \boldsymbol{K}_{\mathrm{G}}\big|} <\tfrac{8}{9};\label{1.24}	
		\end{align}
		\item{Compatibility condition : with  $p_0(x,y)$ be a scalar function}
	\begin{equation}\label{compatibility}
		({\boldsymbol{u}}_{0,h}\cdot\nabla_h){\boldsymbol{u}}_{0,h}
			+\sqrt{\tfrac{\nu}{2 }}
			\nabla_{ h}B\otimes\nabla_{ h}B
			{\boldsymbol u}_{0,h}
		 +\nabla_h p_0=0.
	\end{equation} 
	\end{enumerate}
	Then, for the approximate $\boldsymbol{u}_{app}^\varepsilon$ given in Theorem \ref{thapp}, there exists $T<\min(T^*,\pi/2)$ such that
	\begin{equation}\label{1.25} \lim\limits_{\varepsilon\rightarrow0^+}
		\lVert
		\boldsymbol{u}^\varepsilon
		-\boldsymbol{u}_{app}^\varepsilon 
		\rVert_{{L}^{\infty}([0,T)\times\Omega)}=0\,.
	\end{equation}	
\end{theorem}

		\begin{equation}   
		\lVert{\boldsymbol{u}}_{0,h}\rVert_{{W}^{1,\infty}(\mathbb{R}^2)}\leq\sqrt{\tfrac{\nu}{3}}\,;
		\end{equation}

\begin{lemma}\label{f4}
	If \( g(0) = 0 \), \( g(t) \in C^1 \), and   
	\begin{equation}\label{1.28}
		\frac{d}{dt}g(t) \leq C_1^2 g^2(t) + C_2^2,
	\end{equation}
	where \( C_1, C_2 \) are positive constants, then the inequality 
	\begin{equation}\label{1.29}
		g(t) \leq \tan(C_1 C_2 t) \frac{C_2}{C_1}\,.
	\end{equation}
	holds for time  $t< \frac{\pi}{2C_1 C_2}$.
\end{lemma}

\begin{equation}\label{2.65}
	\begin{cases}
		\partial_{t} {\boldsymbol v}^{\varepsilon}-\nu\varepsilon\Delta{\boldsymbol v}^{\varepsilon}
		+{\boldsymbol u}^{\varepsilon} \cdot \nabla {\boldsymbol v}^{\varepsilon}
		+{\varepsilon}^{-1}\boldsymbol{R}\boldsymbol v^{\varepsilon}
		+{\varepsilon}^{-1}\nabla p_v 
		=-{\boldsymbol v}^{\varepsilon} \cdot \nabla\boldsymbol{u}^\varepsilon_{app}
		-\boldsymbol \rho^{\varepsilon}\,,\\
		\nabla\cdot {\boldsymbol v}^{\varepsilon}=0\,,\\
		{\boldsymbol v}^{\varepsilon}|_{t=0}={\boldsymbol v}^{\varepsilon}|_{\partial\Omega}=0\,.
	\end{cases}
\end{equation}

\begin{proposition}\label{propIC}
	Assuming ${\boldsymbol{u}}_{0,h}\in H^s(\mathbb{R}^2)(s>5)$ and $B(x,y)\in C^\infty(\mathbb{R}^2)$, 
	for $t\in [0,+\infty)$,
	the internal terms $\boldsymbol{u}^I$ and correction terms $\boldsymbol{u}^c$ satisfy 
	\begin{align*} 
		\lVert \nabla\big(
		\boldsymbol{u}^{I} 
		+\boldsymbol{u}^{c}
		\big)(t)
		\rVert_{L^{\infty}(\Omega)}
		\lesssim
		&\Big(\lVert{\boldsymbol{u}}_{0,h}\rVert_{L^2(\mathbb{R}^2)}^2+\lVert{\omega}_0\rVert_{H^2(\mathbb{R}^2)}^2\Big) \mathrm{e}^{-\frac{\sqrt{2\nu}}{8}  t}
		\notag\\
		&+\varepsilon
		\lVert {\boldsymbol{u}}_{0,h} \rVert_{H^{4}(\mathbb{R}^2)}
		+\varepsilon^2
		\big(
		1+
		\lVert {\boldsymbol{u}}_{0,h} \rVert_{H^{4}(\mathbb{R}^2)}
		\big)
		\lVert {\boldsymbol{u}}_{0,h} \rVert_{H^{s}(\mathbb{R}^2)}\,,
	\end{align*}
	and
	\begin{align*} 
		&\lVert \partial_t
		\big(
		\boldsymbol{u}^{I} 
		+\boldsymbol{u}^{c}\big)(t) \rVert_{L^{\infty}(\Omega)}\\
		&\qquad\lesssim
		\lVert{\boldsymbol{u}}_{0,h}\rVert_{H^{3}(\mathbb{R}^2)}
		+\varepsilon
		\lVert{\boldsymbol{u}}_{0,h}\rVert_{H^{4}(\mathbb{R}^2)}
		+\varepsilon^2
		\big(
		1+
		\lVert {\boldsymbol{u}}_{0,h} \rVert_{H^{5}(\mathbb{R}^2)}
		\big)
		\lVert {\boldsymbol{u}}_{0,h} \rVert_{H^{s}(\mathbb{R}^2)}\,.
	\end{align*}
\end{proposition}

\begin{proposition}\label{propBL}
	Assuming $\bar{\boldsymbol{u}}_h\in H^s(\mathbb{R}^2)(s>5)$ and $B(x,y)\in C^\infty(\mathbb{R}^2)$, for $k=0,1$, the boundary term $\boldsymbol{u}^{BL}$  satisfy
	\begin{align*} 
		&\lVert d(z)^{\frac{1}{2}} |\nabla^k\boldsymbol{u}^{BL}|(t) \rVert_{{L}^2([B,B+2];{L}^\infty(\mathbb{R}^2))}
		\\
		&\qquad\lesssim
		\tfrac{3\sqrt{3\nu}}{4}
		\varepsilon^{1-k}
		\big(
		\lVert{\boldsymbol{u}}_{0,h}\rVert_{L^2(\mathbb{R}^2)}+\lVert{\omega}_0\rVert_{H^2(\mathbb{R}^2)}	
		\big)\mathrm{e}^{-\frac{\sqrt{2\nu}}{8}  t}
		\notag\\
		&\qquad\quad+\tfrac{3\sqrt{3\nu}}{4}
		\varepsilon^{1-k}
		\big(
		\varepsilon
		\lVert {\boldsymbol{u}}_{0,h} \rVert_{H^{4}(\mathbb{R}^2)}+
		\varepsilon^2
		\big(
		1+
		\lVert {\boldsymbol{u}}_{0,h} \rVert_{H^{4}(\mathbb{R}^2)}
		\big)
		\lVert {\boldsymbol{u}}_{0,h} \rVert_{H^{s}(\mathbb{R}^2)}
		\big)\,,
	\end{align*}
	and
	\begin{align*} 
		\lVert \partial_t\boldsymbol{u}^{BL} \rVert_{L^\infty(\Omega)}
		\lesssim
		&\lVert{\boldsymbol{u}}_{0,h}\rVert_{H^{3}(\mathbb{R}^2)}
		+\varepsilon
		\lVert{\boldsymbol{u}}_{0,h}\rVert_{H^{4}(\mathbb{R}^2)}\\
		&+\varepsilon^2
		\big(
		1+
		\lVert {\boldsymbol{u}}_{0,h} \rVert_{H^{5}(\mathbb{R}^2)}
		\big)\lVert {\boldsymbol{u}}_{0,h} \rVert_{H^{s}(\mathbb{R}^2)}\,,
	\end{align*}	
	where $d(z)$ denotes the distance to the boundary.
\end{proposition}

\begin{proposition}\label{prop5}
	Under the assumptions of Theorem \ref{thm2}, then there exist some $\varepsilon>0$ such that
	\begin{equation}\label{4.1}
		\lVert\boldsymbol v^{\varepsilon}\rVert^2_{L^\infty([0,T),{L}^2(\Omega))}
		+\nu\varepsilon\lVert\nabla\boldsymbol v^{\varepsilon}\rVert^2_{L^2([0,T),{L}^2(\Omega))}
		\lesssim\varepsilon^4\,.
	\end{equation} 
\end{proposition}

	\begin{equation}\label{4.4}
		|\langle{\boldsymbol u}^{\varepsilon} \cdot \nabla {\boldsymbol v}^{\varepsilon}
		+	{\varepsilon}^{-1}\boldsymbol{R}\boldsymbol v^{\varepsilon}
		+		{\varepsilon}^{-1}\nabla p_v,{\boldsymbol v}^{\varepsilon}\rangle|=0\,.	
	\end{equation}

	\begin{align}\label{4.7}
		&|\langle{\boldsymbol v}^{\varepsilon} \cdot \nabla
		(
		\boldsymbol{u}^I+
		\boldsymbol{u^c}
		),{\boldsymbol v}^{\varepsilon}\rangle|\notag\\
		&\qquad\leq\lVert {\boldsymbol v}^{\varepsilon} \rVert^2_{{L}^2(\Omega)}
		\lVert \nabla
		\big(
		\boldsymbol u^{I}+\boldsymbol u^{c}
		\big)
		\rVert_{{L}^\infty(\Omega)}
		\notag\\
		&\qquad\leq
		\big(
		\lVert{\boldsymbol{u}}_{0,h}\rVert_{L^2(\mathbb{R}^2)}+\lVert{\omega}_0\rVert_{H^2(\mathbb{R}^2)}
		\big)\mathrm{e}^{-\frac{\sqrt{2\nu}}{8}  t}\lVert {\boldsymbol v}^{\varepsilon} \rVert^2_{{L}^2(\Omega)}
		\notag\\
		&\qquad\quad+\big(\varepsilon
		\lVert {\boldsymbol{u}}_{0,h} \rVert_{H^{4}(\mathbb{R}^2)}
		+\varepsilon^2
		\big(
		1+
		\lVert {\boldsymbol{u}}_{0,h} \rVert_{H^{4}(\mathbb{R}^2)}
		\big)
		\lVert {\boldsymbol{u}}_{0,h} \rVert_{H^{6}(\mathbb{R}^2)}\big)\lVert {\boldsymbol v}^{\varepsilon} \rVert^2_{{L}^2(\Omega)}
		\,,
	\end{align}

\begin{proposition}\label{prop8}
	Under the assumptions of Theorem \ref{thm2}, then there exist some $\varepsilon>0$ such that
	\begin{equation}\label{4.1.1}
		\lVert\partial_t\boldsymbol v^{\varepsilon}\rVert^2_{L^\infty([0,T),{L}^2(\Omega))}
		\lesssim\varepsilon^3\,.
	\end{equation}
	In particular, it holds that
	\begin{equation}\label{4.2.1}
		\lVert\nabla\boldsymbol v^{\varepsilon}\rVert^2_{L^\infty([0,T),{L}^2(\Omega))}
		\lesssim\varepsilon^{\frac52}\,.
	\end{equation}
\end{proposition}

\begin{proof}
	First acting the operator $\partial_t$ on the system \eqref{2.65} and then making an inner product with $\partial_t\boldsymbol{v}^\varepsilon$ yields
	\begin{align}
		\label{4.16.1}
		&\frac{1}{2}\frac{d}{dt}\lVert\partial_t\boldsymbol v^{\varepsilon}\rVert^2_{{L}^2(\Omega)}
		+\nu\varepsilon\lVert\nabla\partial_t\boldsymbol v^{\varepsilon}\rVert^2_{{L}^2(\Omega)}\notag\\
		&\qquad\leq
		|\langle{\boldsymbol u}^{\varepsilon} \cdot \nabla \partial_t{\boldsymbol v}^{\varepsilon}
		+	{\varepsilon}^{-1}\boldsymbol{R}
		\partial_t\boldsymbol v^{\varepsilon}
		+		{\varepsilon}^{-1}\nabla \partial_tp_v,\partial_t{\boldsymbol v}^{\varepsilon}\rangle|\notag\\
		&\qquad\quad+
		|\langle
		\partial_t{\boldsymbol v}^{\varepsilon}\cdot
		\nabla\boldsymbol{u}^{\varepsilon}_{app},
		\partial_t{\boldsymbol v}^{\varepsilon}
		\rangle|
		+
		|\langle
		\partial_t{\boldsymbol v}^{\varepsilon} \cdot \nabla\boldsymbol{v}^\varepsilon
		,\partial_t{\boldsymbol v}^{\varepsilon}\rangle|\notag\\
		&\qquad\quad+
		|\langle
		\partial_t{\boldsymbol u}_{app}^{\varepsilon} \cdot \nabla\boldsymbol{v}^\varepsilon
		,\partial_t{\boldsymbol v}^{\varepsilon}\rangle|
		+
		|\langle{\boldsymbol v}^{\varepsilon} \cdot \nabla\partial_t\boldsymbol{u}^\varepsilon_{app},
		\partial_t{\boldsymbol v}^{\varepsilon}\rangle|	
		+|\langle
		\partial_t\boldsymbol \rho^{\varepsilon}	,\partial_t{\boldsymbol v}^{\varepsilon}\rangle|
		\,.
	\end{align}
	
	
	For the first two terms, by applying methods analogous to those used in the $L^2$ estimates of equations \eqref{4.4}-\eqref{4.7}, we can obtain
	\begin{equation}\label{3.2.1}
		|\langle{\boldsymbol u}^{\varepsilon} \cdot \nabla \partial_t{\boldsymbol v}^{\varepsilon}
		+	{\varepsilon}^{-1}\boldsymbol{R}
		\partial_t\boldsymbol v^{\varepsilon}
		+		{\varepsilon}^{-1}\nabla \partial_t p_v,\partial_t{\boldsymbol v}^{\varepsilon}\rangle|=0\,,
	\end{equation}
	and
	\begin{align}\label{3.3.1}
		&|\langle
		\partial_t{\boldsymbol v}^{\varepsilon}\cdot
		\nabla\boldsymbol{u}^{\varepsilon}_{app},
		\partial_t{\boldsymbol v}^{\varepsilon}
		\rangle|\notag\\
		&\qquad\lesssim
		\tfrac{3\sqrt{3\nu}}{4}
		\varepsilon
		\big(
		\lVert{\boldsymbol{u}}_{0,h}\rVert_{L^\infty(\mathbb{R}^2)} 
		+\varepsilon\big(
		\lVert{\boldsymbol{u}}_{0,h}\rVert_{L^2(\mathbb{R}^2)}+\lVert{\omega}_0\rVert_{H^2(\mathbb{R}^2)}
		\big)
		\big)\mathrm{e}^{-\frac{\sqrt{2\nu}}{8}  t}
		\lVert\nabla\partial_t{\boldsymbol v}^{\varepsilon}\rVert^2_{{L}^2(\Omega)}
		\notag\\
		&\qquad\quad+\tfrac{3\sqrt{3\nu}}{4}\varepsilon^3
		\big(
		1+
		\lVert {\boldsymbol{u}}_{0,h} \rVert_{H^{3}(\mathbb{R}^2)}
		\big)
		\lVert {\boldsymbol{u}}_{0,h} \rVert_{H^{5}(\mathbb{R}^2)}
		\lVert\nabla\partial_t{\boldsymbol v}^{\varepsilon}\rVert^2_{{L}^2(\Omega)}\notag\\
		&\qquad\quad+
		\big(
		\lVert{\boldsymbol{u}}_{0,h}\rVert_{L^2(\mathbb{R}^2)}+\lVert{\omega}_0\rVert_{H^2(\mathbb{R}^2)}
		\big)\mathrm{e}^{-\frac{\sqrt{2\nu}}{8}  t}\lVert \partial_t{\boldsymbol v}^{\varepsilon} \rVert^2_{{L}^2(\Omega)}
		\notag\\
		&\qquad\quad+\big(\varepsilon
		\lVert {\boldsymbol{u}}_{0,h} \rVert_{H^{4}(\mathbb{R}^2)}
		+\varepsilon^2
		\big(
		1+
		\lVert {\boldsymbol{u}}_{0,h} \rVert_{H^{4}(\mathbb{R}^2)}
		\big)
		\lVert {\boldsymbol{u}}_{0,h} \rVert_{H^{6}(\mathbb{R}^2)}\big)\lVert\partial_t {\boldsymbol v}^{\varepsilon} \rVert^2_{{L}^2(\Omega)}\,.
	\end{align}
	
	
	For the third term, by employing the Cauchy-Schwarz inequality and the Gagliardo-Nirenberg inequality, we derive the following result
	\begin{align}\label{4.17.1}
		|\langle
		\partial_t{\boldsymbol v}^{\varepsilon} \cdot \nabla\boldsymbol{v}^\varepsilon)
		,\partial_t{\boldsymbol v}^{\varepsilon}\rangle|
		\leq
		&\lVert
		\nabla{\boldsymbol v}^{\varepsilon} 	\rVert_{{L}^2(\Omega)}
		\lVert
		\partial_t{\boldsymbol v}^{\varepsilon} 	\rVert^2_{{L}^4(\Omega)}
		\notag\\
		\leq &
		\lVert\nabla{\boldsymbol v}^{\varepsilon} 	\rVert_{{L}^2(\Omega)}
		\lVert
		\partial_t{\boldsymbol v}^{\varepsilon} 	\rVert^{\frac12}_{{L}^2(\Omega)}
		\lVert
		\nabla\partial_t{\boldsymbol v}^{\varepsilon} 	\rVert^{\frac32}_{{L}^2(\Omega)}
		\notag\\
		\lesssim&
		\frac{\nu\varepsilon}{8}
		\rVert
		\nabla\partial_t{\boldsymbol v}^{\varepsilon} 	\lVert^{2}_{{L}^2(\Omega)}
		+
		\varepsilon^{-3}
		\lVert\nabla{\boldsymbol v}^{\varepsilon} 	\rVert^4_{{L}^2(\Omega)}
		\lVert
		\partial_t{\boldsymbol v}^{\varepsilon} 	\rVert^{2}_{{L}^2(\Omega)}
		\,.
	\end{align}
	Below we estimate item $\lVert\nabla{\boldsymbol v}^{\varepsilon} 	\rVert_{{L}^2(\Omega)}$ in \eqref{4.17.1}. Using Proposition \ref{prop5}, we have  
	\begin{equation*}
		\frac{d}{dt}\lVert\boldsymbol v^{\varepsilon}\rVert^2_{{L}^2(\Omega)}
		+\frac{\nu\varepsilon}{2}
		\lVert\nabla\boldsymbol v^{\varepsilon}\rVert^2_{{L}^2(\Omega)} \lesssim \varepsilon^4 \,,
	\end{equation*}
	which implies  
	\begin{equation*}
		\varepsilon\lVert\nabla\boldsymbol v^{\varepsilon}\rVert^2_{{L}^2(\Omega)} \lesssim \varepsilon^4 
		+\Big|
		\frac{d}{dt}\lVert\boldsymbol v^{\varepsilon}\rVert^2_{{L}^2(\Omega)}
		\Big|\,.
	\end{equation*} 
	Applying the Cauchy-Schwarz inequality, we further derive 
	\begin{align}\label{3.1.1}
		\lVert\nabla\boldsymbol v^{\varepsilon}\rVert^2_{{L}^2(\Omega)} \lesssim
		& \varepsilon^3 
		+\varepsilon^{-1}\Big|
		\int_{\Omega} 2 \boldsymbol v^{\varepsilon} \cdot {\partial_t}\boldsymbol v^{\varepsilon} \, dx
		\Big|\notag\\
		\lesssim&\varepsilon^3 
		+
		2\varepsilon^{-1}\lVert\boldsymbol v^{\varepsilon}\rVert_{{L}^2(\Omega)}
		\lVert\partial_t\boldsymbol v^{\varepsilon}\rVert_{{L}^2(\Omega)}\notag\\
		\lesssim&\varepsilon^3 
		+
		\varepsilon 
		\lVert\partial_t\boldsymbol v^{\varepsilon}\rVert_{{L}^2(\Omega)}\,.
	\end{align} 
	Thus, \eqref{4.17.1} can be rewritten as
	\begin{align}\label{4.17.2}
		|\langle
		\partial_t{\boldsymbol v}^{\varepsilon} \cdot \nabla\boldsymbol{v}^\varepsilon)
		,\partial_t{\boldsymbol v}^{\varepsilon}\rangle|
		\lesssim&
		\frac{\nu\varepsilon}{8}
		\rVert
		\nabla\partial_t{\boldsymbol v}^{\varepsilon} 	\lVert^{2}_{{L}^2(\Omega)}
		+
		\varepsilon^{3}\lVert
		\partial_t{\boldsymbol v}^{\varepsilon} 	\rVert^{2}_{{L}^2(\Omega)}
		+
		\varepsilon^{-1}\lVert
		\partial_t{\boldsymbol v}^{\varepsilon} 	\rVert^{4}_{{L}^2(\Omega)}
		\,.
	\end{align}
	
	
	
	For the nonlinear terms related to the approximate solution, we decompose the approximate solution following the approach in Proposition \ref{prop5}. Using the results from Propositions \ref{propIC} and \ref{propBL} and the incompressible conditions of $\boldsymbol{v}^\varepsilon$ and $\boldsymbol{u}^\varepsilon_{app}$, we obtain
	\begin{align}\label{4.21.1}
		&|\langle
		\partial_t{\boldsymbol u}_{app}^{\varepsilon} \cdot \nabla\boldsymbol{v}^\varepsilon
		,\partial_t{\boldsymbol v}^{\varepsilon}\rangle|
		+
		|\langle{\boldsymbol v}^{\varepsilon} \cdot \nabla\partial_t\boldsymbol{u}^\varepsilon_{app},
		\partial_t{\boldsymbol v}^{\varepsilon}\rangle|	
		\notag\\
		&\qquad=|\langle
		\partial_t{\boldsymbol u}_{app}^{\varepsilon} \cdot \nabla\partial_t\boldsymbol{v}^\varepsilon
		,{\boldsymbol v}^{\varepsilon}\rangle|
		+
		|\langle{\boldsymbol v}^{\varepsilon} \cdot \nabla\partial_t{\boldsymbol v}^{\varepsilon},
		\partial_t\boldsymbol{u}^\varepsilon_{app}\rangle|	
		\notag\\
		&\qquad\leq
		\rVert
		\nabla\partial_t{\boldsymbol v}^{\varepsilon} 	\lVert_{{L}^2(\Omega)}
		\rVert
		{\boldsymbol v}^{\varepsilon} 	\lVert_{{L}^2(\Omega)}
		\rVert
		\partial_t\boldsymbol{u}^\varepsilon_{app}
		\lVert_{{L}^\infty(\Omega)}\notag\\
		&\qquad\lesssim
		\frac{\nu\varepsilon}{8}
		\rVert
		\nabla\partial_t{\boldsymbol v}^{\varepsilon} 	\lVert^{2}_{{L}^2(\Omega)}
		+
		\varepsilon^{-1}\rVert
		{\boldsymbol v}^{\varepsilon} 	\lVert^2_{{L}^2(\Omega)}
		\rVert
		\partial_t\boldsymbol{u}^\varepsilon_{app}\lVert^2_{{L}^\infty(\Omega)}
		\notag\\
		&\qquad\lesssim
		\frac{\nu\varepsilon}{8}
		\rVert
		\nabla\partial_t{\boldsymbol v}^{\varepsilon} 	\lVert^{2}_{{L}^2(\Omega)}
		+
		\varepsilon^3
		\,.
	\end{align}
	
	
	
	
	According to Theorem \ref{thapp} and , the last term is rearranged as
	\begin{align}\label{4.19.1}
		|\langle
		\partial_t\boldsymbol \rho^{\varepsilon}	,\partial_t{\boldsymbol v}^{\varepsilon}\rangle|
		\leq&
		\lVert
		\partial_t\boldsymbol{\rho}^\varepsilon\rVert_{{L}^{2}(\Omega)}
		\rVert
		\partial_t{\boldsymbol v}^{\varepsilon}
		\lVert_{{L}^2(\Omega)}
		\notag\\	
		\lesssim&
		\rVert
		\partial_t{\boldsymbol v}^{\varepsilon}
		\lVert^2_{{L}^2(\Omega)}
		+\varepsilon^4
		\lVert
		{\boldsymbol{u}}_{0,h}
		\rVert^2_{{H}^{6}(\mathbb{R}^2)}
		\big(
		\lVert
		{\boldsymbol{u}}_{0,h}
		\rVert^2_{{H}^{4}(\mathbb{R}^2)}
		+1
		\big)^2
		\,.
	\end{align}

	
	
	
	According to \eqref{3.2.1}-\eqref{4.19.1}, we can rewrite equation \eqref{4.16.1} as follows
	\begin{align}
		\label{4.20.1}
		&\frac{1}{2}\frac{d}{dt}\lVert\partial_t\boldsymbol v^{\varepsilon}\rVert^2_{{L}^2(\Omega)}
		+\frac{3\nu\varepsilon}{4}\lVert\nabla\partial_t\boldsymbol v^{\varepsilon}\rVert^2_{{L}^2(\Omega)}\notag\\
		&\qquad\lesssim
		\tfrac{3\sqrt{3\nu}}{4}
		\varepsilon
		\big(
		\lVert{\boldsymbol{u}}_{0,h}\rVert_{L^\infty(\mathbb{R}^2)} 
			+\varepsilon\big(
		\lVert{\boldsymbol{u}}_{0,h}\rVert_{L^2(\mathbb{R}^2)}+\lVert{\omega}_0\rVert_{H^2(\mathbb{R}^2)}
		\big)
		\big)\mathrm{e}^{-\frac{\sqrt{2\nu}}{8}  t}
		\lVert\nabla\partial_t{\boldsymbol v}^{\varepsilon}\rVert^2_{{L}^2(\Omega)}
		\notag\\
		&\qquad\quad+
		\tfrac{3\sqrt{3\nu}}{4}\varepsilon^3
		\big(
		1+
		\lVert {\boldsymbol{u}}_{0,h} \rVert_{H^{3}(\mathbb{R}^2)}
		\big)
		\lVert {\boldsymbol{u}}_{0,h} \rVert_{H^{5}(\mathbb{R}^2)}
		\lVert\nabla\partial_t{\boldsymbol v}^{\varepsilon}\rVert^2_{{L}^2(\Omega)}\notag\\
		&\qquad\quad
		+
		\varepsilon^{-1}\lVert
		\partial_t{\boldsymbol v}^{\varepsilon} 	\rVert^{4}_{{L}^2(\Omega)}
		+		(1+\varepsilon^{3})\lVert
		\partial_t{\boldsymbol v}^{\varepsilon} 	\rVert^{2}_{{L}^2(\Omega)}+
		\varepsilon^3
		\,.
	\end{align}
	
	
	
	When  $\varepsilon$ is sufficiently small, we can simplify equation \eqref{4.20.1} with the following expressions:
	\begin{align}
		 \varepsilon^{-1} \lVert
		\partial_t{\boldsymbol v}^{\varepsilon} 	\rVert^{4}_{{L}^2(\Omega)}&\leq
		\frac{
			\lVert
			\partial_t{\boldsymbol v}^{\varepsilon} 	\rVert^{4}_{{L}^2(\Omega)}
		}{
			\varepsilon^3 
		}\,,\label{1}\\
		(1+\varepsilon^{3} )\lVert
		\partial_t{\boldsymbol v}^{\varepsilon} 	\rVert^{2}_{{L}^2(\Omega)}&\leq
		\frac{
			\lVert
			\partial_t{\boldsymbol v}^{\varepsilon} 	\rVert^{4}_{{L}^2(\Omega)}
		}{
			\varepsilon^3 
		}
		+\varepsilon^3 \,.\label{2}
	\end{align}
	
	
	Under condition  \eqref{1.26} in Theorem \ref{thm2}, if $\varepsilon$ is sufficiently small, we obtain 
	\begin{equation}
		\label{4.20.2}
		\frac{1}{2}\frac{d}{dt}\lVert\partial_t\boldsymbol v^{\varepsilon}\rVert^2_{{L}^2(\Omega)}
		\lesssim
		\frac{
			\lVert
			\partial_t{\boldsymbol v}^{\varepsilon} 	\rVert^{4}_{{L}^2(\Omega)}
		}{
			\varepsilon^3
		}
		+\varepsilon^3\,.
	\end{equation}
	
	Finally, by applying Lemma \ref{f4} for $t< \frac\pi2$, we derive the result \eqref{4.1.1}, and subsequently obtain \eqref{4.2.1}.
\end{proof}

\end{document}
